% --- INICIO DE PLANTILLA PERSONALIZADA ---
\documentclass[oneside,11pt]{Latex/Classes/PhDthesisPSnPDF}

% --- Paquetes Originales ---
\usepackage{tabularx}
\usepackage{float}
\usepackage{multirow, array}
\usepackage[usenames,dvipsnames,svgnames,table]{xcolor}
\usepackage{longtable}
\usepackage{latexsym,amsmath,amssymb,amsfonts}
\usepackage{enumitem}
\usepackage{xparse}
\usepackage{booktabs}
\usepackage[round, sort, numbers]{natbib}  
\usepackage{adjustbox}
\usepackage[utf8]{inputenc}
\usepackage{caption}
\usepackage{graphicx}
\usepackage{hyperref}
\usepackage{cleveref}

% Definición de colores
\definecolor{miblue}{RGB}{68, 114, 196}

% --- CAMBIO CRÍTICO 1: Usar \input para macros en el preámbulo ---
% \include da problemas antes del begin document. \input es seguro.
\input{Latex/Macros/MacroFile1}

% --- Definiciones de seguridad (por si la clase falla al cargar alguna variable) ---
\providecommand{\facultad}[1]{\def\lafacultad{#1}}
\providecommand{\escudofacultad}[1]{\def\elescudo{#1}}
\providecommand{\degree}[1]{\def\elgrado{#1}}
\providecommand{\director}[1]{\def\eldirector{#1}}
\providecommand{\lugar}[1]{\def\ellugar{#1}}
\providecommand{\degreedate}[1]{\def\lafecha{#1}}

% --- Configuración para bloques de código de R (Pandoc) ---

%%%%%%%%%%%%%%%%%%%%%%%%%%%%%%%%%%%%%%%%%%%%%%%%%%%%%%%%%%%%%%%%%%%%%%%%%%%%%%%%
%                         DATOS (Conectados a RMarkdown)                       %
%%%%%%%%%%%%%%%%%%%%%%%%%%%%%%%%%%%%%%%%%%%%%%%%%%%%%%%%%%%%%%%%%%%%%%%%%%%%%%%%
\title{El Ranking ATP como Indicador de Éxito en el Tenis de Élite
(2016-2025)}
\author{Andrés Ortega \\ Marco Ramirez}

% Lógica para Facultad: Si la defines en el YAML la usa, si no, usa la default

\facultad{Facultad de Ciencias Económicas y Sociales \\ Escuela de
Estadistica y Ciencias Actuariales}
  \escudofacultad{Latex/Classes/Escudos/faces}

\degree{Computación I}
\director{Jesus Ochoa \\ Oliver Triveño}
\degreedate{Abril 2026} 
\lugar{Caracas}

\portadatrue 

% Metadatos PDF
\keywords{}
\subject{}

% --- CORRECCIÓN PARA PANDOC NUEVO (Imágenes) ---
\makeatletter
\providecommand{\pandocbounded}[1]{#1}
\makeatother


%%%%%%%%%%%%%%%%%%%%%%%%%%%%%%%%%%%%%%%%%%%%%%%%%%%%%
%                   DOCUMENTO                       %
%%%%%%%%%%%%%%%%%%%%%%%%%%%%%%%%%%%%%%%%%%%%%%%%%%%%%
\begin{document}

\maketitle

%%%%%%%%%%%%%%%%%%%%%%%%%%%%%%%%%%%%%%%%%%%%%%%%%%%%%
%                  PRÓLOGO                          %
%%%%%%%%%%%%%%%%%%%%%%%%%%%%%%%%%%%%%%%%%%%%%%%%%%%%%
\frontmatter

% Puedes agregar dedicatorias desde el YAML si quieres, o dejarlas fijas aquí

%%%%%%%%%%%%%%%%%%%%%%%%%%%%%%%%%%%%%%%%%%%%%%%%%%%%%
%                   ÍNDICES                         %
%%%%%%%%%%%%%%%%%%%%%%%%%%%%%%%%%%%%%%%%%%%%%%%%%%%%%
\setcounter{secnumdepth}{3}
\setcounter{tocdepth}{3}

\tableofcontents



\cleardoublepage

  \cleardoublepage       % Empieza en página derecha (estándar en tesis)
  \chapter*{Resumen}     % Crea el título "Resumen" sin numeración (Capítulo *)
  \addcontentsline{toc}{chapter}{Resumen} % (Opcional) Lo añade al índice
  
  En el presente informe se realizará un análisis exploratorio de datos
basado en un histórico de 10 temporadas entre 2016-2025 del Top 50 en el
circuito ATP del tenis profesional, el Ranking ATP es la métrica clave
para evaluar el desempeño profesional de un jugador, pero, existen otros
variables como el tipo de superficie, la categoría de la competición y
el avance de rondas en un torneo, que hacen que esta confiable métrica
no siempre sea tan representativa. De acuerdo con lo anteriormente
expuesto, se examinará por medio de estadísticas descriptivas como el
ranking ATP fluctúa dependiendo del contexto en el que se desarrolle la
competencia.             % Aquí se pega el texto que escribiste en el YAML
  
  \cleardoublepage

%%%%%%%%%%%%%%%%%%%%%%%%%%%%%%%%%%%%%%%%%%%%%%%%%%%%%
%                   CONTENIDO (El Corazón)          %
%%%%%%%%%%%%%%%%%%%%%%%%%%%%%%%%%%%%%%%%%%%%%%%%%%%%%
\mainmatter
\def\baselinestretch{1.5}

% --- CAMBIO CRÍTICO 2: Aquí RMarkdown inyecta todo tu texto ---
\chapter*{Introducción}

En el área del deporte profesional, la recopilación y análisis de datos
históricos es una tarea fundamental para generar estadísticas
descriptivas que permitan recopilar información relevante del
rendimiento de un jugador o equipo. A diferencia de los deportes jugados
en equipo dónde las estadísticas pueden llegar a atenuarse por el
desempeño grupal, en el caso del tenis, debido a su estructura de
competición individual (single), cada dato (puntos obtenidos, ranking,
cantidad de sets jugados) manifiesta el desempeño individual de cada
atleta. A menudo, se asume que el ranking de un jugador en un partido es
una métrica clave la cual expresa cual es la tendencia de quien pueda
llegar a ser el ganador, sin embargo, existen otras variables (ronda del
torneo, superficie de la cancha, tipo de serie) las cuales pueden llegar
a influir en el desempeño de jugadores y dejan a un lado al Ranking ATP
como factor determinante.\\

Por lo tanto, este estudio busca responder la siguiente pregunta: ¿Es el
ranking de un jugador una métrica suficiente para reflejar lo que
sucederá en cancha?

\chapter{Planteamiento del Problema}

En el tenis profesional, el ranking ATP es la métrica oficial que mide
el éxito y el desempeño de los jugadores, tal como se explica en la guía
técnica de Wilson (sf) el ranking ATP funciona como un barómetro de
consistencia, donde cada torneo tiene un peso específico en la
puntuación total del atleta. Este sistema no solo determina quiénes son
los mejores del mundo, sino que también condiciona la participación en
torneos, la asignación de cabezas de serie y las expectativas del
público. Ser parte del top 50 no solo garantiza notoriedad en el tenis,
sino que te ofrece otros beneficios como el acceso a torneos de mayor
categoría, una exposición constante a diferentes superficies, mayores
incentivos financieros y prestigio en la carrera deportiva, Association
of Tennis Profesionals (ATP, 2026).\\

Ahora bien, ¿Es el ranking ATP una métrica suficiente para comparar a
dos jugadores que compiten en cancha? En el presente trabajo
describiremos cómo se comportan los jugadores del top 50 en diferentes
superficies, rondas y torneos en un histórico comprendido entre
2016-2025 en la época dorada del tenis.\\

Con respecto a lo anteriormente expuesto, se plantean las siguientes
preguntas de investigación:\\

\textbf{1.} ¿Cómo se comporta el ranking de los ganadores y perdedores
en el ATP tour?\\

\textbf{2.} ¿En cúales superficies el ranking ATP es más
representativo?\\

\textbf{3.} ¿En qué torneos se enfocan los jugadores con rankings más
altos y más bajos?\\

\textbf{4.} ¿Cúal es el ranking promedio predominante en las finales de
los torneos más importantes?\\

\section{Justificación}

La justificación del presente trabajo se sustenta en evaluar la robustez
del Ranking ATP como un descriptor del desempeño real en el tenis
profesional.\\

En el ámbito deportivo, el ranking de un jugador es una métrica
acumulativa de resultados pasados \hspace{0pt}\hspace{0pt}que definen su
estatus en el circuito. Esta investigación busca que por medio de la
estadística descriptiva identificar los patrones de distribución de los
ganadores y perdedores. Para ello, se tomará como universo en estudio a
los partidos de 10 temporadas consecutivas en el período entre 2016-2025
de ambos atletas pertenecientes al top 50, este segmentación se realiza
producto a que estos atletas tienen acceso directo a los Grand Slam,
Masters 1000 y a una gran parte de los torneos ATP 500 y ATP 250 . Por
otra parte, tienen prohibida la participación en eventos de menor
categoría sin una carta de invitación (Association of Tennis
Professionals, 2026, sec.~7.07), esto garantizá que el enfoque de este
estudio se base en el circuito de élite del tenis.\\

\section{Objetivos}

\subsection{Objetivo General}
\begin{itemize}
  \item{Determinar si existe una relación entre el posicionamiento en el ranking ATP y el desempeño competitivo de los jugadores para evaluar su rendimiento en los partidos según superficie, torneo y ronda en los atletas pertenecientes al Top 50 durante el periodo 2016-2025.}
\end{itemize}

\subsection{Objetivos Específicos}
\begin{itemize}
  \item{Comparar el desempeño de los ganadores y perdedores de los jugadores pertenecientes al top 50 en el período de análisis.}
  \item{Analizar si la posición del ranking de los jugadores varía de manera sustancial de acuerdo a la superficie del partido jugado.}
  \item{Verificar la eficiencia del ranking como pronosticador del éxito mediante el avance en las rondas.}
  \item{Identificar si los jugadores con mejor posicionamiento (Top 10), concentran su participación en eventos con mayor prestigio.}
\end{itemize}

\section{Cobertura}

\subsection{Horizontal}

La cobertura horizontal se basa en cada uno de los partidos de tenis
registrados en la base de datos. Se considera el momento, lugar y torneo
en el cual se efectuaron los partidos, y el resultado de los mismos
junto a ciertas características que describen los enfrentamientos.\\

\begin{table}[ht]
\centering
\caption{Variables consideradas en el estudio} 
\label{tab:variables}
\renewcommand{\arraystretch}{1.4}%
\begin{tabularx}{\textwidth}{llX} %
  \hline
\textbf{TIPO} &\textbf{VARIABLE} &\textbf{DESCRIPCIÓN}\\ 
  \hline
Categórica & Tournament & Nombre oficial del torneo dónde se disputo el encuentro\\ 
Numérica & Date & Fecha en la que se llevo a cabo el partido(Formato AAAA-MM-DD)\\
Categórica & Series & Categoría o nivel del torneo\\
Categórica & Surface & Tipo de suelo de la cancha: Duro (Hard), Arcilla (Clay) o Hierba (Grass)\\
Categórica & Round & Etapa eliminatoria del torneo\\
Categórica & Player\_1 & Nombre del primer jugador registrado en el encuentro\\
Categórica & Player\_2 & Nombre del segundo jugador registrado en el encuentro\\
Categórica & Winner & Nombre del jugador que resultó victorioso en el encuentro\\
Numérica & Rank\_1 & Posición en el ranking ATP del Jugador 1 al iniciar el torneo\\
Numérica & Rank\_2 & Posición en el ranking ATP del Jugador 2 al iniciar el torneo\\
Numérica & Winnner\_Rank & Posición en el ranking ATP del ganador del partido\\
Numérica & Loser\_Rank & Posición en el Ranking ATP del perdedor del partido\\
Numérica & Tasa de Victorias & Porcentaje de victorias anuales de cada jugador por temporada\\ 
   \hline
\end{tabularx}
\end{table}

\subsection{Vertical}

La cobertura vertical del trabajo de investigación comprende un total de
66790 observaciones. Cada fila representa un partido único jugado en el
ATP tour. El alcance temporal de los datos en longitudinal, abarcando
desde el año 2000 hasta principios del 2026, lo que permite un análisis
del desempeño de jugadores de diferentes épocas y con distintos estilos
de juego.

\section{Periodo de Referencia}

El presente trabajo de investigación se sustenta en una base de datos
históricos del ATP tour, que abarca una serie de registros de más de dos
décadas. El periodo de referencia definido para el análisis comprende
desde el 1 de enero de 2000 hasta el primer trimestre de 2026.

\section{Variables}

Las dos variables principales en el estudio (Variables
dependientes/objetivo) son Rank\_1 y Rank\_2, que representa el ranking
ATP de un cada jugador en ese juego específico, como variables
explicativas o independientes se analizan:

\begin{enumerate}
  \item Variables temporales: \textbf{Date}
  \item Variables categóricas: \textbf{Tournament}, \textbf{Series}, \textbf{Surface}, \textbf{Round},             \textbf{Jugador\_1}, \textbf {Jugador\_2}, \textbf{Winner}
  \item Variables numericas: \textbf{Winner\_Rank}, \textbf{Loser\_Rank}, \textbf{Tasa de victorias} 
\end{enumerate}

\chapter{Marco Teórico}

\section{Bases Teóricas}

El presente capítulo establece los fundamentos matemáticos y
estadísticos necesarios para el análisis de los datos del historial de
partidos del ATP Tour. De igual forma, se abordan las definiciones del
sistema de clasificación, la categorización de los torneos y los
factores externos que influyen en la consistencia de los resultados
deportivos.\\

\subsection{El sistema de clasificación de la ATP}

El ranking de la Asociación de Tenistas Profesionales (ATP) es la
herramienta objetiva utilizada para determinar la jerarquía de los
jugadores y su acceso a los torneos. Se basa en un sistema de méritos
acumulativos durante un período de 52 semanas, el cual suma los mejores
19 resultados de un jugador determinado en dicho período (Wilson, sf).

\begin{itemize}
\item \textbf{Puntos de mérito:}Los jugadores acumulan puntos según la ronda alcanzada y la importancia del torneo.
\end{itemize}

\subsection{Categorización de torneos y estructura de puntos}

El circuito profesional se divide en distintas Series, las cuales
definen la cantidad de puntos en disputa y el nivel de competitividad.\\

\textbf{1.} Grand Slams: Los cuatro torneos de mayor prestigio (Abierto
de Australia, Roland Garros, Wimbledon y US Open), que otorgan 2000
puntos al ganador.\\

\textbf{2.} ATP Masters 1000: Nueve torneos anuales que representan el
segundo nivel de importancia.\\

\textbf {3.} ATP 500 y 250: Torneos de menor escala que sirven para la
rotación y acumulación de puntos de jugadores en diversos rangos del
ranking.\\

\subsection{ Factores del entorno en el rendimiento competitivo}

\begin{itemize}
\item
  Superficie (Surface): Es el tipo de pista en la cual se lleva a cabo
  el juego, hay 3 tipos de pista: Arcilla (Clay), Césped (Grass) y Pista
  Dura (Hard). Cada una de ellas determina una dinámica distinta de
  juego ya que cambian la altura, y velocidad de la pelota, así como el
  deslizamiento del jugador. La estadística deportiva demuestra que en
  el tenis existen jugadores ``especialistas'' cuyo rendimiento varía
  significativamente según el suelo, independientemente de su ranking
  general.\\
\item
  Rondas (Round): A medida que un torneo progresa, el nivel de
  resistencia del oponente aumenta. Estadísticamente, se espera que la
  desviación estándar del ranking de los ganadores merme conforme se
  avanza hacia las rondas finales.\\
\end{itemize}

\subsection{Análisis exploratorio de datos}

En virtud de que la presente investigación tiene un enfoque descriptivo,
un Análisis Exploratorio de datos (EDA) permite identificar patrones,
detectar valores atípicos y resumir las características principales del
estudio.\\

Matemáticamente Sea\(X\) una variable que representa el ranking del
ganador e \(Y\) una variable que representa el ranking del perdedor,
entonces:\\

\begin{itemize}
\item \textbf{Media muestral}
\end{itemize}

\[\bar{x}=\frac{\sum_{i=1}^{n} x_i}{n} \vee_{}^{} \bar{y}=\frac{\sum_{i=1}^{n} y_i}{n}\]

\begin{itemize}
 \item \textbf{Varianza(desvíos con respecto a la media)}: Indica la desviación promedio del ranking de los ganadores o perdedores con respecto a su media.
\end{itemize}

\[S^2_x=\frac{\sum_{i=1}^{n}(x_i- \bar{x})^2} {n-1} \vee S^2_y=\frac{\sum_{i=1}^{n}(y_i- \bar{y} )^2} {n-1}\]
\newpage

\begin{itemize}
  \item \textbf{Tasa de victorias}
\end{itemize}

\[T.V= \frac {P.G}{P.J}*100\%\] Donde,

\begin{itemize}
\item P.G representa al número de partidos en los que el jugador resultó ganador en una temporada.
\item P.J representa el numéro total de partidos que disputo el jugador en la temporada.
\end{itemize}

\hfill\break

\begin{itemize}
\item \textbf{Coeficiente de correlación de Pearson}
\end{itemize}

\[\rho_{X,Y} = \frac{Cov(X,Y)}{\sigma_X \sigma_Y} = \frac{E[(X - \mu_X)(Y - \mu_Y)]}{\sigma_X \sigma_Y}\]
En este estudio, tomaremos como la variable \(X\) al ranking promedio de
un jugador en una temporada e \(Y\) representará la tasa de victorias de
dicho jugador en la temporada.

El coeficiente de correlación de Pearson permitirá determinar si el
ranking de un jugador es un descriptor del éxito elcual,a priori, a
medida de que aumenta su ranking, se espera que también lo haga su tasa
de victorias.

\chapter{Marco Metódico}

En esta sección se presentan las metodologías y/o test necesarios
relacionados con la práctica divulgada en el marco teórico, así como las
fuentes de datos.

\section{Bases metódicas utilizadas en el trabajo de investigación}

En esta sección se presentan las metodologías

\subsection{Item 1}

Descripción\\

\subsection{Item 2}

Descripción\\

\subsubsection{Item 2.2}

Descripción\\

\subsection{Item n}

Descripción\\

\chapter{Análisis de Resultados}

\{\small Este capítulo presenta los resultados obtenidos al aplicar las
metodologías descritas en el capítulo anterior, así como del análisis
estadístico descriptivo de las variables estudiadas y las ventajas y
desventajas al aplicar dichos tests.\}

\section{Presentación de resultados 1}

Descripción\\

\section{Presentación de resultados 2}

Descripción\\

\section{Presentación de resultados n}

Descripción\\

\chapter{Conclusiones y Recomendaciones}

\textbf{1.-} Conclusion 1 del trabajo de investigación.\\

\textbf{Recomendación} Recomendación 1 del trabajo de investigación.\\

\textbf{2.-} Conclusion 2 del trabajo de investigación.\\

\textbf{Recomendación} Recomendación 2 del trabajo de investigación.\\

\textbf{n.-} Conclusion n del trabajo de investigación.\\

\textbf{Recomendación} Recomendación n del trabajo de investigación.\\

\appendix
\renewcommand{\thechapter}{\Alph{chapter}}

\chapter{Anexo A}

\chapter{Anexo B}

\chapter{Anexo C}

%%%%%%%%%%%%%%%%%%%%%%%%%%%%%%%%%%%%%%%%%%%%%%%%%%%%%
%                   REFERENCIAS                     %
%%%%%%%%%%%%%%%%%%%%%%%%%%%%%%%%%%%%%%%%%%%%%%%%%%%%%
% Mantenemos la estructura natbib de tu plantilla
\renewcommand{\bibname}{Lista de Referencias}
\bibliographystyle{apalike} 
\bibliography{bibliografia.bib} 

%%%%%%%%%%%%%%%%%%%%%%%%%%%%%%%%%%%%%%%%%%%%%%%%%%%%%
%                   APÉNDICES                       %
%%%%%%%%%%%%%%%%%%%%%%%%%%%%%%%%%%%%%%%%%%%%%%%%%%%%%

\end{document}
